% 张三慧《大学物理学》第4版V1《力学、热学》，第7章《波动》
% 原教材扫描PDF物理页248--253，印刷232--237页。
% 本稿保留习题7.1--7.33的源编号与数学条件；思考题7.x不混入。
% 图7.44--7.51、图7.47地震时间坐标与原书公式(7.45)仍须独立图文校验。
\section*{第7章\quad 波动}
\subsection*{章末习题7.1--7.33}

\paragraph{习题 7.1}
太平洋一次海浪速度为
\(740\,\mathrm{km/h}\)，波长
\(300\,\mathrm{km}\)。
求波频率及波横越
\(8000\,\mathrm{km}\) 海域所用时间。

\paragraph{习题 7.2}
横波沿一条长弦以
\(u=0.80\,\mathrm{m/s}\) 传播。
在 \(x=0.10\,\mathrm m\) 处测得
\[
 y=0.05\sin(1.0-4.0t)
\]
（SI单位）。写出整条弦的行波表达式。
\textit{原题未指定传播方向，正、负 \(x\) 两个分支都应讨论。}

\paragraph{习题 7.3}
已知沿绳传播横波
\[
 y=2.0\times10^{-2}
 \sin[2\pi(200t-2.0x)].
\]
求：（1）波长、频率、传播速率及方向；
（2）绳上质元振动最大速率与波速之比。

\paragraph{习题 7.4}
报道1976年唐山地震时地面曾向上抬起
约 \(2\,\mathrm m\)。
把地震波简化为频率
\(1\,\mathrm{Hz}\)、传播速率
\(3\,\mathrm{km/s}\) 的简谐波。
求波长与地面质点的最大振动速率。

\paragraph{习题 7.5}
图7.44为向 \(+x\) 传播的平面横波在
\(t=0\) 时的波形。
图上振幅 \(0.040\,\mathrm m\)，
相邻向上过零点相距 \(0.40\,\mathrm m\)，
\(x=0\) 处向上过零，波速
\(u=0.080\,\mathrm{m/s}\)。
（1）求波函数；（2）画
\(t=T/8\) 时的波形。
\textit{图示来源：V1 PDF物理页248、图7.44。}

\paragraph{习题 7.6}
某波满足
\[
 y=A\cos[\pi(4t+2x)].
\]
（1）求 \(t=4.2\,\mathrm s\)
的位移—位置表达式，并求此刻距原点最近
的一个波峰位置与它经过原点的时刻；
（2）画 \(t=4.2\,\mathrm s\) 的波形。

\paragraph{习题 7.7}
简谐波频率
\(500\,\mathrm{Hz}\)，波速
\(350\,\mathrm{m/s}\)。
（1）沿传播方向相差 \(60^\circ\)
的两质点相距多远？
（2）同一点相隔
\(10^{-3}\,\mathrm s\) 的两振动状态相差多少相位？

\paragraph{习题 7.8}
钢内纵波传播速率
\(5100\,\mathrm{m/s}\)，
密度
\(7.8\times10^3\,\mathrm{kg/m^3}\)。
在简化一维弹性介质模型中求钢的杨氏模量。

\paragraph{习题 7.9}
设均匀弹性固体或液体
体积弹性模量为 \(K\)，
微小体积形变为 \(\Delta V\)，
原体积为 \(V\)。
证明体积形变的弹性势能密度
\[
 u_{\rm elastic}=\frac12K
 \left(\frac{\Delta V}{V}\right)^2,
 \qquad|\Delta V|\ll V.
\]

\paragraph{习题 7.10}
纵波在钢中以
\(u=5.1\times10^3\,\mathrm{m/s}\) 传播，
振幅 \(A=10^{-9}\,\mathrm m\)，
频率 \(f=10^3\,\mathrm{Hz}\)，
钢密度
\(7.9\times10^3\,\mathrm{kg/m^3}\)，
钢轨横截面积
\(S=15\,\mathrm{cm^2}\)。
求波强（平均能流密度）和波传输的平均功率。

\paragraph{习题 7.11}
参照教材图7.12，
取简谐行波传播中一段长
\(\Delta x\) 的绳元，
从绳张力对截面位移速度所作功出发，
证明行波的平均能流表达式
\(P_{\rm av}=\tfrac12\mu\omega^2A^2u\)，
其中 \(\mu\) 为绳线密度；
说明其与教材波强公式(7.31)一致。

\paragraph{习题 7.12}
A、B两点有相干波源，振幅相等，
频率同为
\(100\,\mathrm{Hz}\)，初相相差
\(\pi\)。
源间距
\(AB=30\,\mathrm m\)，
波速 \(400\,\mathrm{m/s}\)。
求两源之间AB连线上振动叠加后
增强且相对静止的点的位置。

\paragraph{习题 7.13}
驻波表达式
\[
 y=0.02\cos(20x)\cos(750t).
\]
求：（1）形成它的两行波振幅与速率；
（2）相邻波节间距；
（3）\(t=2.0\times10^{-3}\,\mathrm s\)，
\(x=5.0\times10^{-2}\,\mathrm m\)
处介质质点振动速度。

\paragraph{习题 7.14}
正向传播的平面简谐波振幅
\(A\)、频率 \(\nu\)、波速 \(u\)，
图7.45中反射面位于距原点
\(3\lambda/4\) 的 \(P\) 处。
（1）\(t=0\) 时原点处质元位于平衡位置
并沿 \(+y\) 运动，写入射波函数；
（2）波由波疏介质传播到波密反射面，
反射振幅与入射相等，写反射波函数，
求正向入射与左向反射波的节点。
\textit{图7.45：反射面相对于原点位于右侧。}

\paragraph{习题 7.15}
石英薄片通过上下金属电极施加交流电产生
厚度方向振动，片厚
\(d=2.00\,\mathrm{mm}\)；
石英厚度方向纵波声速
\(u=5.74\times10^3\,\mathrm{m/s}\)。
以厚度方向基频驻波的
\(\lambda=2d\) 条件估算共振电压频率。

\paragraph{习题 7.16}
一项吉尼斯纪录对应声级
\(115\,\mathrm{dB}\)，另一纪录为
\(141\,\mathrm{dB}\)。
以标准参考声强
\(I_0=10^{-12}\,\mathrm{W/m^2}\)，
求对应的声强数值。

\paragraph{习题 7.17}
图7.47地震波记录显示
P波与S波到达时刻约相差
\(5\,\mathrm{min}\)。
其平均传播速率分别为
\(8\,\mathrm{km/s}\) 与
\(6\,\mathrm{km/s}\)。
估算震中到台站的距离。
\textit{源图时标必须独立复核，不能把分钟误当作秒。}

\paragraph{习题 7.18}
1976年唐山地震按题设震级
\(M=9.2\) 估计，
利用教材式(7.45)的震级—能量经验关系
计算总释放能。
若作比较，求相当于多少吨TNT的能量；
原书按1kg TNT约
\(4.6\times10^6\,\mathrm J\)。

\paragraph{习题 7.19}
停泊海岸的船随海浪上下振动，
周期
\(4.0\,\mathrm s\)，
振幅 \(0.60\,\mathrm m\)；
连续浪峰相隔
\(25\,\mathrm m\)。
求海波传播速率以及海面上介质质点绕圆运动
的线速度量级，比较这两种速度。

\paragraph{习题 7.20}
摩托车驾驶员以
\(80\,\mathrm{km/h}\) 行驶，
后方警车为追赶它以
\(120\,\mathrm{km/h}\) 前进，
警笛本征频率 \(400\,\mathrm{Hz}\)，
空气声速
\(330\,\mathrm{m/s}\)。
求摩托车驾驶员听到的警笛频率。

\paragraph{习题 7.21}
海面波长
\(120\,\mathrm m\)、周期
\(10\,\mathrm s\)。
快艇速度
\(24\,\mathrm{m/s}\)。
求快艇迎浪（逆波传播）与顺浪行驶时
所遇浪峰的频率和间隔时间。

\paragraph{习题 7.22}
一海豚在水下向接近的潜艇发送
\(1.8\times10^4\,\mathrm{Hz}\)
超声脉冲，
潜艇反射的回声与海豚原发声频率相差
\(220\,\mathrm{Hz}\)。
水中声速
\(1.54\times10^3\,\mathrm{m/s}\)。
在潜艇近似静止的条件下估算海豚游泳速度。

\paragraph{习题 7.23}
血液流速约
\(0.32\,\mathrm{m/s}\)。
超声波沿血流方向发射
\(4.0\,\mathrm{MHz}\)，
被红细胞反射后原超声接收器接收。
求回波频率与发射频率的差，
取人体组织声速
\(1.54\times10^3\,\mathrm{m/s}\)。
说明血流和声波方向对频移符号的影响。

\paragraph{习题 7.24}
高速公路测速雷达频率
\(5.0\times10^{10}\,\mathrm{Hz}\)，
反射波与入射波形成拍频
\(1.1\times10^4\,\mathrm{Hz}\)。
估计迎面汽车速度并判断
是否超过限速
\(100\,\mathrm{km/h}\)。

\paragraph{习题 7.25}
飞机在
\(5.00\times10^3\,\mathrm m\) 高度
以马赫数 \(2.3\) 水平飞行，
声速 \(330\,\mathrm{m/s}\)。
求：（1）马赫锥半顶角；
（2）飞机飞过地面观察者头顶后
经过多长时间，观察者才听到冲击波。

\paragraph{习题 7.26}
船以
\(60\,\mathrm{km/h}\)
在千岛湖水面行驶，身后船迹波的张角为
\(10^\circ\)（按图7.49示意）。
用船波尾迹形成的波锥几何关系估算
静水波的传播速率。
\textit{需核实图7.49的10°是半角还是全角；
本稿采用源文字「张角10°」与图意的条件计算。}

\paragraph{习题 7.27}
两个同方向波叠加：
\[
 y_1=A\sin(5x-10t),\qquad
 y_2=A\sin(4x-9t).
\]
求：（1）合成波表达式；
（2）波群的速度；
（3）一个波包（包络单个波瓣）的长度，
并说明是否应将包络的符号正负误认为波能为负。

\paragraph{习题 7.28}
沿均匀细棒的弯曲波频散关系
\[
 \omega=ak^2,\qquad a>0.
\]
求波相速度与群速度之间的关系。

\paragraph{习题 7.29}
电离层中无线电波的频散关系
\[
 \omega^2=\omega_p^2+c^2k^2,
\]
其中 \(\omega_p\) 为常数、
\(c\) 为真空光速。
求相速度 \(u_p\) 与群速度
\(u_g\)，证明
\(u_pu_g=c^2\)。

\paragraph{习题 7.30}
某遥远星系发出的光
经过测量，其光波长比地球上同类原子光谱
波长增长 \(3/2\) 倍。
在把红移解释为\emph{狭义相对论径向多普勒}
的模型中，求远离地球的速率；
注意经典非相对论近似不能在超光速结果上继续使用。

\paragraph{习题 7.31}
证明教材图7.38所示的有限波包
长度与谱宽关系
\[
 \Delta x=\frac{\lambda^2}{\Delta\lambda},
 \qquad
 \Delta x\Delta k=2\pi,
\]
并以波包延续时间
\(\Delta t=\Delta x/u_g\) 说明
\(\Delta t\Delta\nu=1\) 的近似关系。
声明这些符号的\emph{谱宽及包络定义}，
避免把精确的傅里叶不等式误当成所有波形的严格等式。

\paragraph{习题 7.32}
已知光在反射/折射介质中沿
总传播时间为极值（最短驻时）的路径传播。
依图7.50(a)(b)，设入射光源、
反射或折射点、接收点分别在界面两侧给定，
通过最小传播时间关于界面交点的变化，
推导反射定律
\(i'=i\) 以及折射定律
\(\sin i/\sin r=u_1/u_2\)。

\paragraph{习题 7.33}
超声电机原理见图7.51。
薄金属片上各质元同时参与互相垂直的两振动：
\[
 \xi_y=A_y\sin(\omega t-kx),\qquad
 \xi_x=A_x\cos(\omega t-kx).
\]
（1）证明质元合运动轨迹是
\[
 \frac{\xi_x^2}{A_x^2}+
 \frac{\xi_y^2}{A_y^2}=1.
\]
（2）证明金属片与滑块接触处的水平运动速度
\[
 v_x=-\omega A_x\sin(\omega t-kx)
\]
及向左运动的符号约定。
